\section{Conclusiones}
\label{sec:conclusiones}

Se especificó y validó formalmente, mediante VDM++, el componente
crítico de una plataforma de planificación de viajes en transporte
público para Arequipa. El modelo define los tipos del dominio, las
operaciones de búsqueda de rutas directas y combinadas, y sus
postcondiciones, que formalizan las garantías del algoritmo.

La validación en la VDM++ Toolbox permitió comprobar que el modelo es
sintáctica y semánticamente consistente, que las operaciones se ejecutan
correctamente sobre casos de prueba representativos y que las
postcondiciones se satisfacen. De este modo, se evidenció que aplicar
métodos formales a un algoritmo de movilidad permite especificar con
precisión su comportamiento y detectar ambigüedades que las pruebas
informales no revelan.
